\subsection{Zero-mode contributions to the determinant of the unit-diagonal matrix}\label{sec:B4T08}
Retain the notation of Sections~\ref{sec:B4T06} and~\ref{sec:B4T07}. Write $y=p^*_{\rm hijo}-w$ and
\[
 f_i=\sum_{\{a,b\}\subset a_i}\Gamma_{ab},\qquad
 f_i^I=\sum_{\{a,b\}\subset a_i\cap I}\Gamma_{ab},
 \qquad K_0=\Gamma-I.
\]
The sums run over unordered pairs of distinct lines. Each $f_i$ is $O(k^{-1})$.
The perturbation $\Delta_0=S_{0,\rm hijo}-S_{0,II}$ has zero diagonal and satisfies
\[
 \operatorname{tr}(\Gamma_{II}\Delta_0)
 =-2\sum_Fp_i f_i^I+2\sum_{E'}y_i f_i.
\]
The identity is exact because each pair of lines intersects in at most one square.

\paragraph{Positive Schur complement and covolume.}
To handle the zero mode without an ill-conditioned gauge-fixed inverse, set $\Pi=vv^T/(2k)$ and $\widetilde S_0=S_0+\Pi$.
The latter has inverse $\Gamma+\Pi$. Let
\[
 C_J=S_{0,JJ}+\Pi_{JJ},\quad V_J=S_{0,IJ}+\Pi_{IJ},\quad
 P_*=\widetilde S_{0,II}-V_J C_J^{-1}V_J^T,
\]
\[
 H_*=P_*^{-1}=(\Gamma+\Pi)_{II},\quad
 \Pi'=v_Iv_I^T/[2(k-1)],\quad
 D_*=\Delta_0+(\Pi'-\Pi_{II})+V_JC_J^{-1}V_J^T.
\]
The exact Schur and cofactor identities give
\[
 \frac{\det S_{0,\rm hijo,gf}}{\det S_{0,\rm padre,gf}}
 =\frac{k}{k-1}\frac{\det(I+H_*D_*)}{\det C_J}.
\]
The factor $k/(k-1)$ appears only when these regularized matrices are introduced. It is consistent with the direct cofactor identity of Section~\ref{sec:B4T01}.

\paragraph{Four first-order calculations.}
Since $S_{0,JJ}=I+p_c(\mathbf1\mathbf1^T-I)$,
\[
 \log[k/(k-1)]-\log\det C_J=O(k^{-2}).
\]
The two projector corrections satisfy
\[
 \operatorname{tr}(\Pi_{II}\Delta_0)=-1/k+O(k^{-2}),\qquad
 \operatorname{tr}\big[H_*(\Pi'-\Pi_{II})\big]=1/k+O(k^{-2}).
\]
Here we use $\Delta_0v_I=S_{0,IJ}v_J$ and $\Gamma v=0$.
For the rank-four term, block inversion gives
\[
 \operatorname{tr}(H_*V_JC_J^{-1}V_J^T)
 =\operatorname{tr}\big[C_J(\Gamma+\Pi)_{JJ}-I\big]
 =\sum_{j\in J}(K_0)_{jj}+2/k+O(k^{-2}).
\]

\paragraph{Contribution of the quadratic term.}
The leading part of $\Delta_0$ is the sparse matrix with entry $-p_i$ for each pair among the three retained lines of a square $i\in F\setminus c$. Its $\infty$-norm is $O(k^{-1})$. The remaining entries of $D_*$ are $O(k^{-2})$, with absolute row sums $O(k^{-1})$.
Since $H_*=I+O_{\max}(k^{-1})$ and $\|H_*\|_\infty=O(1)$,
\[
 \operatorname{tr}(H_*D_*H_*D_*)=6\sum_Fp_i^2+O(k^{-2}).
\]
Distinct squares give disjoint sets of line pairs. The contribution of $c$ is only $O(k^{-2})$. The cross terms involving the dense part are $O(k^{-2})$, as is the cubic remainder in the logarithm of the determinant. Combining these calculations gives
\[
 \boxed{\log\frac{\det S_{0,\rm hijo,gf}}{\det S_{0,\rm padre,gf}}
 =\sum_{j\in J}(K_0)_{jj}+\frac2k+
 \operatorname{tr}(\Gamma_{II}\Delta_0)-3\sum_Fp_i^2+O(k^{-2}).}
\]

\paragraph{Identity for first-order cancellation.}
Summing the four diagonal entries of $S_0\Gamma=I-\Pi$ indexed by the removed lines gives
\[
 \boxed{\sum_Fp_i(f_i-f_i^I)
 =-\sum_{j\in J}(K_0)_{jj}-\frac1k-p_cf_c.}
\]
Each square in $F\setminus c$ meets exactly one line in $J$. The square $c$ counts each pair twice in the resolvent instead of once. The exact difference is $p_cf_c=O(k^{-2})$.
This identity uses the exact zero mode and requires no independence assumption.
